% =============================================================================
%  Conclusion générale
%
%  Aucun chiffre n'est introduit ici qui ne figure déjà, avec son contexte et
%  sa marge d'erreur, au chapitre 5. La conclusion synthétise ; elle ne
%  renforce pas.
% =============================================================================

\chapter*{Conclusion générale}
\addcontentsline{toc}{chapter}{Conclusion générale}
\label{chap:conclusion}
\markboth{CONCLUSION GÉNÉRALE}{CONCLUSION GÉNÉRALE}

% =============================================================================
\section*{Ce qui était demandé, ce qui a été livré}
% =============================================================================

L'énoncé d'EURAFRIC Information demandait un système d'optimisation de
portefeuille piloté par apprentissage automatique, capable d'apprendre les
corrélations dynamiques entre actifs, opérant sous des contraintes réalistes de
gestion, et évalué dans un cadre rigoureux de \emph{backtesting} hors
échantillon.

Le livrable est une chaîne complète et reproductible~: collecte et validation de
données provenant de cinq sources hétérogènes~; une couche de préparation en
trois niveaux, validée à l'écriture~; un moteur de simulation temporelle dont
l'absence de vision du futur est vérifiée par des tests plutôt qu'affirmée~;
quatre familles de modèles ajoutées une à une, du rétrécissement de covariance à
la détection de régime, de la prédiction de rendement à l'optimisation sensible
aux coûts~; un protocole d'évaluation comportant validation croisée purgée,
segment de test gelé, intervalles de confiance et correction pour le nombre
d'essais~; et deux surfaces d'usage — une application web à deux pages et une
interface REST — partageant une couche de données unique. L'ensemble est couvert
par \textbf{390~tests} exécutés hors ligne en moins de deux minutes.

% =============================================================================
\section*{Ce que le projet établit}
% =============================================================================

Les résultats sont classés selon leur force probante, et cette distinction est
elle-même un résultat de la démarche.

\paragraph{Démontré.} Le détecteur de régime, non supervisé et n'ayant jamais
reçu la moindre liste de crises, a identifié les cinq crises étudiées~: 91,7~\%
de ses décisions sont classées « baissier » à l'intérieur des fenêtres de crise
contre 29,2~\% en dehors, soit un rapport de \textbf{3,13}, avec un dépassement
du taux de base sur \textbf{cinq crises sur cinq} ($p = 0{,}031$ au test des
signes). C'est le seul résultat statistiquement significatif du projet. Sa portée
doit être énoncée avec précision~: le modèle \emph{voit} ce qu'il prétend voir, en
temps réel et sans connaître la suite~; le fait que le système \emph{gagne} à agir
sur cette information est, lui, indémontrable à cette taille d'échantillon.

\paragraph{Observé, non démontré.} Le système complet dépasse la meilleure
approche classique de \textbf{6,2~\%} de ratio de Sharpe net sur le portefeuille
visé, et lui est inférieur de \textbf{1,6~\%} sur l'univers ETF. Les intervalles
de confiance se chevauchent largement, le classement change selon le protocole
d'évaluation retenu, et le gain provient d'une réduction de volatilité et non
d'un supplément de rendement. Ces trois réserves sont affichées dans le rapport,
dans le tableau de bord et dans l'interface de programmation — cette dernière
livrant la réserve de manière indissociable du chiffre.

\paragraph{Observé, robuste.} Pendant les cinq crises, l'optimisation sous
contrainte a systématiquement mieux protégé le capital que la diversification
naïve, avec un délai de retour au sommet environ deux fois plus court.
L'attribution est faite honnêtement~: ce gain revient à la contrainte de position
et au modèle de covariance, non à la couche de régime, qui dans trois de ces cinq
fenêtres produit exactement la même allocation qu'une stratégie classique.

\paragraph{Découvert.} Le plafond de 25~\% par actif, retenu comme règle métier et
non comme dispositif de modélisation, s'est révélé le régulariseur le plus
efficace du système~: le relâcher coûte 10,1~\% de ratio de Sharpe, davantage que
l'écart entre deux modèles quelconques testés sur cet univers. Ce résultat est la
reproduction, sur nos propres données, d'un théorème connu — une contrainte de
position active équivaut mathématiquement à un rétrécissement de la matrice de
covariance.

\paragraph{Réfuté.} Cinq pistes indépendantes visant à dépasser la ligne de base
ont été explorées et fermées~: la prédiction de rendement calibrée honnêtement,
un historique cinq fois plus long, l'ajout de données fondamentales, la
re-sélection à chaque repli temporel, et le conditionnement de la contrainte au
régime. Le constat qui émerge de leur convergence est plus intéressant que chaque
échec pris isolément~: \textbf{à cette taille d'échantillon, l'évaluation ne peut
pas résoudre des écarts de l'ordre de grandeur que ces idées produisent}. Deux de
ces études établissent en outre, séparément, que \emph{précision de prédiction et
performance de portefeuille sont deux objectifs distincts}~: le coefficient
d'information a doublé sans que le ratio de Sharpe suive, et il a même reculé
dans un cas.

% =============================================================================
\section*{Ce que le projet apporte au-delà des chiffres}
% =============================================================================

La contribution la plus durable de ce travail n'est probablement pas une
performance, mais un \emph{dispositif}.

Un incident l'illustre mieux qu'un discours. À trois semaines de la fin, la
découverte que les dividendes marocains manquaient dans les données a réduit de
plus de moitié le gain publié du système. Le chiffre corrigé a remplacé l'ancien
partout — rapport, tableau de bord, livrables — et deux affirmations défendues
depuis le début du projet ont dû être retirées. Cette correction n'a été possible
que parce que la chaîne était reconstructible de bout en bout à partir des
sources~; elle n'aurait probablement pas été détectée du tout dans un projet où
les résultats sont produits à la main.

De là vient le principe qui structure l'ingénierie du projet~: \textbf{une règle
importante n'est pas documentée, elle est rendue exécutable}. L'impossibilité de
consulter le futur, l'interdiction d'écrire un chiffre de performance en dur dans
une interface, la cohérence entre documents publiés, la traçabilité de chaque
fonction vers le problème qu'elle traite, l'hermétisme réseau de la suite de
tests~: chacune de ces règles est un test qui échoue si on l'enfreint. Les trois
incidents les plus coûteux du projet ont tous consisté en une règle respectée par
discipline jusqu'au jour où elle ne l'a plus été, sans que rien ne le signale.
Chacun a donné lieu à un garde-fou automatique, validé en rejouant l'incident
réel.

% =============================================================================
\section*{Limites}
% =============================================================================

Elles sont énoncées ici plutôt que laissées à découvrir.

\begin{itemize}
  \item \textbf{La taille d'échantillon est la limite dominante.} Le portefeuille
        visé ne compte qu'environ cinq années d'historique exploitable. C'est ce
        qui empêche de conclure, et aucun modèle supplémentaire ne le corrigera.
  \item \textbf{L'historique marocain gratuit ne remonte pas avant 2021}, ce qui
        exclut la crise de 2020 de l'univers principal. La conception à double
        univers en atténue l'effet~; elle ne le supprime pas.
  \item \textbf{Un rendement de dividende sur quatre a dû être estimé}, la source
        publique ne le publiant pas. L'analyse de sensibilité correspondante est
        documentée, mais la donnée reste une estimation.
  \item \textbf{La correction pour le nombre d'essais est bornée} au périmètre de
        chaque recherche, et non à l'ensemble des expériences menées au cours du
        projet — définition volontairement explicite, mais moins sévère que
        l'idéal.
  \item \textbf{Le système n'a pas été exposé à des conditions d'exploitation
        réelles}~: ni exécution d'ordres, ni flux temps réel, ni couverture du
        risque de change dirham/dollar, exclus du périmètre dès le cadrage.
\end{itemize}

% =============================================================================
\section*{Perspectives}
% =============================================================================

Quatre suites se dégagent, classées par rapport valeur/effort.

\textbf{Étendre l'historique marocain.} C'est la seule action qui attaque la
limite dominante. Un historique long a été obtenu manuellement pour dix-sept
valeurs de la Bourse de Casablanca lors d'une des études~; l'industrialiser —
raccorder cet historique profond aux données récentes sur leur période de
recouvrement, puis l'intégrer à la chaîne — donnerait au portefeuille visé la
profondeur dont il manque, y compris la crise de 2008.

\textbf{Adopter le protocole imbriqué comme protocole de référence.} Il a réduit
la largeur des intervalles de confiance de 28,6~\% pour un coût de calcul modéré,
et il a montré que les réglages sélectionnés changent d'un repli à l'autre — ce
qui confirme qu'une valeur unique et globale est inadaptée.

\textbf{Calibrer la pénalité de rotation par modèle.} Les mesures ont établi
qu'une valeur globale unique dégrade le bon modèle tout en corrigeant le mauvais.
La rendre spécifique à chaque modèle est un travail de calibration circonscrit,
au protocole déjà disponible.

\textbf{Adoucir le basculement de régime.} Le basculement franc actuel produit un
pic de transactions à chaque frontière de régime. Un mélange continu, pondéré par
la probabilité de régime, est une alternative documentée dès la conception et non
encore mesurée.

% =============================================================================
\section*{Bilan}
% =============================================================================

Ce projet aura moins appris à l'équipe comment faire gagner un modèle qu'à
\emph{savoir si un modèle gagne}. C'est un apprentissage inconfortable~: il
consiste, l'essentiel du temps, à construire les instruments qui feront tomber
les résultats auxquels on tenait — un test qui interdit de voir le futur, un
intervalle de confiance qui absorbe l'écart qu'on voulait annoncer, une donnée
corrigée qui divise un gain par deux à trois semaines de la remise.

C'est pourtant, dans un domaine où les performances publiées sont notoirement
difficiles à reproduire, la seule compétence qui se transporte d'un sujet à
l'autre. Le système livré à EURAFRIC Information ne prétend pas battre les
méthodes classiques~: il dit précisément dans quelles conditions il les dépasse,
de combien, avec quelle incertitude, et dans quels cas il leur est inférieur. Un
outil qui sait dire cela est plus utile à un gestionnaire qu'un outil qui promet
davantage.
